\scanpage{69}
\deleted{\unsure{Le $Q$-hom|L'élément|Rappelons}} le $Q$-hom \deleted{On écrit le diagramme \unsure{com}}
\[
u_Z:B_X\longrightarrow A_Y\qquad(\text{mod isog})
\]
défini par $Z$ en vertu de 2.1. n'est autre que $u_D^{-1}$.

\uline{Dém}. On utilise le diagramme commutatif
\diagram{\begin{tikzcd}[ampersand replacement=\&,column sep=large,row sep=large]
H^{2n-1}(X,Q_l)(n-1)\arrow[r,"{a(\gamma(D))}"]\arrow[d,"\unsure{\wr|2}\,\deleted{\unsure{\text{isom}}}","\psi_X"'] \& H^1(Y,Q_l)\\
H^1(B_X,\mathbb Q_l)\arrow[r,"u_D^{(1)}"] \& H^1(A_Y,Q_l)\arrow[u,swap,"\wr\quad\varphi_Y"]
\end{tikzcd}}
\unsure{qui|Qui|Ceci} prouve que $a(\gamma(D))$ est un isom. ssi $u_D^{(1)}$ l'est, i.e. $u_D$ est une isogénie, Pour conclure à l'existence de $Z$, il faut montrer que $\varphi_Y^{-1}$ et $\psi_X^{-1}$ sont $Q$-algébriques (de façon indépendante de $\unsure{l|k}$). \deleted{Pour ceci, il faut \unsure{pr}}\deleted{\inserted{On le sait déjà}} Pour $\psi_X$, on est donc ramené à prouver 2.7. dans le cas particulier de l'hom can.
\[
H^{2n-1}(X,Q_l)(n-1)\longrightarrow H^1(B_X,Q_l)
\]
défini par $D\in\operatorname{Cl}(X,B_X)$, pour trouver un élément can. $Z\in A_1(X\times B_X)$.
% archive sheet number: 68

